% --- Chapter 6
\chapter{Kết luận và Hướng phát triển}
\label{ch:conclusion}

Đề tài được thực hiện trong bối cảnh cuộc cải cách hành chính toàn quốc năm 2025 tạo ra ``cú sốc dữ liệu'' chưa có tiền lệ với hạ tầng thông tin doanh nghiệp Việt Nam. Khung giải pháp \textbf{VN~Address~Intelligence} (VNAI) được thiết kế và hiện thực với tham vọng giải quyết đồng thời bốn bài toán mà chưa nghiên cứu nào trước đây giải quyết một cách tích hợp: đồng bộ tự động danh mục hành chính, chuẩn hoá địa chỉ phi cấu trúc bằng AI đa tầng, xác thực không gian bằng thuật toán hình học, và làm giàu dữ liệu đa nguồn theo mô hình thác nước. Chương này tổng kết các kết quả đạt được, đối chiếu với mục tiêu nghiên cứu, làm rõ đóng góp khoa học, ghi nhận các hạn chế còn tồn tại, và đề xuất các hướng phát triển ưu tiên.

% =====================================================================
\section{Tổng kết các kết quả đạt được}
\label{sec:concl-summary}

Các kết quả của đề tài có thể được phân thành ba nhóm theo phương diện trình bày: kiến trúc và triển khai, định lượng đo được, và phương pháp khoa học.

\subsection{Kết quả về kiến trúc và triển khai}
\label{sec:concl-arch}

Đề tài đã hiện thực hoá một nền tảng thống nhất bao gồm sáu thành phần cốt lõi. \emph{Thứ nhất}, cơ sở dữ liệu quan hệ PostgreSQL với phân tách bốn schema chuyên biệt (\schema{mat}, \schema{osm}, \schema{ath}, \schema{prq}) làm nền tảng dữ liệu bền vững cho master hành chính, dữ liệu OpenStreetMap, hub AI và hàng đợi xử lý. \emph{Thứ hai}, lớp API REST trên FastAPI với hơn sáu mươi \thuat{endpoint} cùng giao diện web tĩnh phục vụ vận hành và thực nghiệm. \emph{Thứ ba}, luồng đồng bộ tự động từ nguồn danh mục hành chính NSO kết hợp cơ chế SCD Type 2 \cite{kimball2013} và đồ thị quan hệ đơn vị (\macode{mat.unit\_edge}) cho phép quản lý vòng đời thực thể hành chính với các quan hệ \macode{MERGES\_INTO}, \macode{SPLIT\_FROM}, \macode{RENAMES\_TO}. \emph{Thứ tư}, pipeline AI xếp tầng tích hợp ba mô hình chuyên biệt: PhoBERT cho NER \cite{nguyen2020phobert}, mGTE cho retrieval \cite{mgte2024}, và họ Qwen cho tinh chỉnh \cite{QwenTeam2024} với lượng tử hoá 4-bit \cite{Dettmers2022BNB}. \emph{Thứ năm}, module geospatial xử lý polygon, kiểm tra điểm thuộc vùng, báo cáo \thuat{mismatch} và \thuat{edge inject} cho hiệu chỉnh GPS gần biên. \emph{Thứ sáu}, khung đánh giá SUPA-Bench trên ground truth chỉ đọc với khả năng tái lập đầy đủ.

\subsection{Kết quả định lượng từ artifact}
\label{sec:concl-quantitative}

Các phép đo đã thực hiện trên artifact cho bốn nhóm kết quả định lượng chính.

\emph{Một, chất lượng mô hình NER có giám sát.} PhoBERT NER đạt F1 kiểm định (seqeval) \(\approx \VNAIGENNERFOnePct\,\%\) và token accuracy \(\approx \VNAIGENNERTokenAccPct\,\%\) trên tập \VNAIGENNEREvalN\, mẫu kiểm định công khai --- vượt mức tham khảo của các nghiên cứu trước về NER địa chỉ tiếng Việt nhưng chưa đạt ngưỡng nội bộ Gate~B \(96\%\) đồng thời cho hai chỉ số.

\emph{Hai, chất lượng dữ liệu lineage.} Audit bridge cho thấy phần lớn bản ghi hàng đợi khớp lineage và phân cấp master ở mức trên \(96\%\), với quy mô \VNAIGENAuditQueueTotal\, dòng trong lần đo --- một quy mô đủ tin cậy về mặt thống kê và phản ánh tải nghiệp vụ thực tế đáng kể.

\emph{Ba, chất lượng end-to-end với kịch bản oracle.} SUPA-Bench đã chứng minh chuỗi đánh giá end-to-end hoạt động đúng (\VNASUPAEMvTwoPct\,\% với \VNASUPANRequested\, mẫu trong lần chạy đơn và \(100{,}00\,\% \pm 0{,}00\) với năm lần chạy phân tầng K~=~5, mỗi lần \(2.000\) mẫu), đồng thời làm rõ giới hạn diễn giải khi so với chuỗi tham chiếu tiền cải cách v1 (EM@v1 \(\approx 16{,}14\,\% \pm 0{,}44\)). Bảng tổng hợp K~=~5 đã được \thuat{persist} thành một dòng trong \macode{ath.supa\_stratified\_eval\_summary} với phạm vi \macode{run\_id} 56--60 và \macode{git\_commit} đầy đủ, tạo \emph{bằng chứng tái lập} cho mọi vòng thẩm định trong tương lai.

\emph{Bốn, chất lượng end-to-end với pipeline thật --- kết quả trọng yếu của đề tài.} Nghiên cứu cắt bỏ thành phần (\thuat{ablation study}) trên Google Colab GPU T4 với \VNAIABLATIONTotalSpecimens\, specimens (\VNAIABLATIONNumConfigs\, cấu hình \(\times\) \VNAIABLATIONNPerConfig\, mẫu), commit \texttt{\VNAIABLATIONGitCommit}, đã cung cấp các con số định lượng chính thức cho kiến trúc lai ghép. Cấu hình A1\_FULL (NER + mGTE + LLM) đạt EM@v2 \(= \VNAIABLATIONAOneEMvTwoPct\,\%\), F1 Đường \(= \VNAIABLATIONAOneFOneStreetPct\,\%\), F1 Phường \(= \VNAIABLATIONAOneFOneWardPct\,\%\), F1 Quận \(= \VNAIABLATIONAOneFOneDistrictPct\,\%\) và độ trễ trung bình \VNAIABLATIONAOneLatencyMs\, ms/mẫu --- vượt hầu hết các ngưỡng kỳ vọng đặt ra ban đầu. Cấu hình so sánh A2\_NER\_TFIDF, A2\_NER\_MGTE và A3\_MGTE\_ONLY cùng đạt EM@v2 \(= \VNAIABLATIONATwoEMvTwoPct\,\%\), trong khi A4\_NER\_LLM (không retrieval) chỉ đạt \VNAIABLATIONAFourEMvTwoPct\,\% --- chứng minh \emph{retrieval là thành phần then chốt không thể bỏ qua}, và LLM chỉ phát huy giá trị khi được tích hợp đúng trong kiến trúc lai ghép cùng retrieval (đóng góp định lượng \(+\VNAIABLATIONLLMContributionPp\) điểm phần trăm EM).

\subsection{Kết quả về phương pháp khoa học}
\label{sec:concl-method}

Đề tài cố định bốn thực hành phương pháp khoa học làm chuẩn mực cho mọi thực nghiệm: (i) \thuat{provenance} đầy đủ qua bốn yếu tố (\macode{git\_commit}, \macode{rng\_seed}, \macode{noise\_profile\_id}, \macode{source\_note}); (ii) tách bạch đo lường giữa chất lượng dữ liệu lineage (audit) và chất lượng mô hình AI (NER, E2E); (iii) phân biệt rõ kịch bản oracle với kịch bản pipeline thật để tránh phóng đại năng lực; (iv) lưu trữ bảng tổng hợp trong schema \schema{ath} với payload JSON đầy đủ, cho phép truy hồi và đối chiếu sau này.

% =====================================================================
\section{Đối chiếu với Mục tiêu Nghiên cứu}
\label{sec:concl-objectives}

Ba câu hỏi nghiên cứu được đặt ra ở Mục~\ref{sec:research-questions} nay đối chiếu với kết quả đạt được như sau.

\textbf{RQ1 --- Mô hình hoá biến động hành chính theo thời gian.} Đề tài đã thiết lập đầy đủ mô hình SCD Type 2 trên schema \schema{mat} với cặp \macode{valid\_from}/\macode{valid\_to}, các cờ \macode{is\_active}, \macode{is\_deleted}, cột \macode{predecessor\_id} và \macode{version\_id}. Đồ thị \macode{mat.unit\_edge} với bốn kiểu quan hệ định lượng mọi biến đổi địa giới. API tra cứu lịch sử theo thời điểm cho phép truy vấn trạng thái đơn vị bất kỳ thời điểm nào. \emph{Kết luận:} RQ1 đã được giải quyết hoàn chỉnh ở tầng dữ liệu và API.

\textbf{RQ2 --- Hybrid PhoBERT + mGTE so với tìm kiếm từ vựng truyền thống.} Đề tài đã hiện thực pipeline HYBRID\_V1 tám bước với ba tầng AI tích hợp. Trên artifact hiện tại, PhoBERT NER đạt F1 \(\approx \VNAIGENNERFOnePct\,\%\) --- vượt mức tham khảo trong các nghiên cứu trước tại Việt Nam (Chương~\ref{ch:literature-review}). Nghiên cứu cắt bỏ thành phần trên \VNAIABLATIONTotalSpecimens\, specimens (Mục~\ref{sec:exp-ablation}) cung cấp bằng chứng định lượng \emph{trực tiếp} cho câu hỏi này: cấu hình A2\_NER\_TFIDF (retrieval từ vựng truyền thống) và A2\_NER\_MGTE (retrieval ngữ nghĩa) cùng đạt EM@v2 \(= \VNAIABLATIONATwoEMvTwoPct\,\%\) trên cohort phổ thông, cho thấy lợi thế của mGTE so với baseline lexical chưa thể hiện rõ ở cấp độ chuỗi hoàn chỉnh; trong khi đó, cấu hình lai ghép đầy đủ A1\_FULL bổ sung tầng LLM đạt \VNAIABLATIONAOneEMvTwoPct\,\%, cao hơn cả hai phương án retrieval thuần \VNAIABLATIONLLMContributionPp\, điểm phần trăm. \emph{Kết luận:} RQ2 đã có bằng chứng định lượng cho thấy giá trị thực sự của hybrid không nằm ở việc thay TF-IDF bằng embedding ngữ nghĩa, mà nằm ở tầng tinh chỉnh LLM trên top-\(k\) ứng viên retrieval. Việc đánh giá lợi thế của mGTE trên các cohort khó hơn (D2 nhiễu cao, D3 lưỡng thời) còn là khoảng trống đo lường cho các vòng thực nghiệm tiếp.

\textbf{RQ3 --- Thuật toán hình học không gian cho tự hiệu chỉnh.} Đề tài đã đề xuất và hiện thực ba chiến lược hiệu chỉnh polygon (Buffer-Union, Concave Hull, Edge Injection) với cơ sở lý thuyết tại Chương~\ref{ch:co-so-ly-thuyet}. Endpoint \macode{POST /api/spatial/subdivide} đã vận hành điểm thuộc vùng dựa trên PostGIS với cơ chế \thuat{fallback} sang \macode{ST\_Distance} khi điểm sát biên. \emph{Kết luận:} RQ3 đã có kiến trúc đầy đủ; đánh giá định lượng hiệu quả của các chiến lược cần dữ liệu giao nhận thực tế quy mô lớn.

\textbf{Đối chiếu với Gate~B nội bộ.} F1 và token accuracy NER trên artifact hiện tại \emph{chưa} vượt \(96\%\) đồng thời như Gate~B yêu cầu. Điều này định vị rõ ``khoảng cách còn lại'' giữa bằng chứng khái niệm (\thuat{proof-of-concept}) và mục tiêu vận hành nghiêm ngặt. Đề tài chủ trương báo cáo trung thực khoảng cách này thay vì che giấu, vì nó là cơ sở cho công tác cải tiến tiếp theo.

\textbf{Hạng mục báo cáo định lượng còn chưa hoàn tất.} (i) Bảng E2E theo ba kịch bản S1--S3 trên địa chỉ thực tế. (ii) Cột retrieval (R@k, MRR) trong tổng hợp SUPA-Bench --- cần chạy \macode{evaluate\_retriever.py} với \macode{--persist-db} trên snapshot mô hình đã chọn. (iii) Số liệu pipeline thật trên cohort phân tầng K~=~5 (N~=~2.000 mỗi lần) --- ablation hiện tại (Mục~\ref{sec:exp-ablation}) đã chứng minh kiến trúc đạt ngưỡng EM trên cohort phổ thông; bước tiếp theo là chạy A1\_FULL trên cohort phân tầng D1--D4 để đánh giá độ bền trên các tầng nhiễu nặng. (iv) SUPA-Bench Final với \(N = 10.000\) trên pipeline A1\_FULL nhằm khẳng định mức EM ổn định trên quy mô gấp đôi ablation hiện tại. (v) Phân tích P95 latency và throughput đo từ pipeline đối với A1\_FULL trên hạ tầng production.

% =====================================================================
\section{Đóng góp Khoa học của Đề tài}
\label{sec:concl-contributions}

Đề tài đóng góp vào tri thức khoa học trên ba phương diện: lý luận, phương pháp luận và thực tiễn.

\subsection{Đóng góp lý luận}
\label{sec:concl-theoretical}

\emph{Thứ nhất}, đề tài đề xuất \textbf{mô hình ``Chuẩn hoá địa chỉ nhận thức thời gian'' (Temporal-Aware Address Standardization)} thông qua tích hợp SCD Type 2 \cite{kimball2013} với đồ thị quan hệ \macode{unit\_edge}. Khác với các nghiên cứu trước coi dữ liệu hành chính là tĩnh, mô hình này cho phép pipeline AI xử lý đồng thời cả hai phiên bản hành chính (tiền và hậu cải cách 2025) mà không gặp hiện tượng \thuat{catastrophic forgetting}. Đây là khoảng trống tri thức số một được xác định ở Mục~\ref{sec:research-gaps} và nay đã có một giải pháp hoàn chỉnh.

\emph{Thứ hai}, đề tài mở rộng mô hình \textbf{Address Confidence Score (ACS)} với bốn thành phần chuẩn hoá theo Weighted Sum Model \cite{hwang1981, saaty1980, zavadskas2019multiple}:
\begin{equation*}
\mathrm{ACS}(a_i \mid q) = \alpha \cdot S_{\text{text}} + \beta \cdot S_{\text{sem}} + \gamma \cdot V_{\text{hier}} + \delta \cdot V_{\text{temporal}}
\end{equation*}
trong đó việc đưa \(V_{\text{hier}}\) (kiểm tra phân cấp) và \(V_{\text{temporal}}\) (trọng số thời gian) làm thành phần cứng là điểm mới so với các công thức ACS chỉ dùng độ khớp văn bản và ngữ nghĩa. Bốn trạng thái quyết định (Auto-Accept, Auto-Convert, Suggest, Reject) phản ánh trực tiếp bối cảnh dữ liệu lưỡng thời.

\emph{Thứ ba}, đề tài đề xuất \textbf{ba chiến lược hiệu chỉnh đa giác không gian} (Buffer-Union, Concave Hull, Edge Injection) học từ lịch sử toạ độ giao nhận thực tế. Đây là một nhánh mở rộng cho cơ sở lý thuyết về Hệ thống Thông tin Địa lý (GIS) trong bối cảnh ngày càng nhiều doanh nghiệp logistics tích luỹ dữ liệu định vị thực tế quy mô lớn.

\subsection{Đóng góp phương pháp luận}
\label{sec:concl-methodological}

\emph{Thứ nhất}, đề tài thiết lập \textbf{khung đánh giá SUPA-Bench} có thể tái lập đầy đủ trên ground truth chỉ đọc. Bất biến nghiên cứu cốt lõi --- pipeline SUPA chỉ \macode{SELECT} chứ tuyệt đối không \macode{INSERT} / \macode{UPDATE} / \macode{DELETE} trên \macode{prq.ground\_truth} --- đảm bảo tách bạch giữa nguồn tham chiếu và quy trình đánh giá. Đây là một đóng góp phương pháp luận có thể chuyển giao cho các nghiên cứu tương tự trong miền dữ liệu nghiệp vụ.

\emph{Thứ hai}, đề tài đề xuất \textbf{hệ thống chỉ số P-* và S-*} tách bạch chất lượng tham số mô hình (P-F1, P-Acc) khỏi chất lượng cấu trúc đầu ra (S-NER-EM, S-E2E-EM, S-RET). Cách phân loại này giúp cô lập nguyên nhân khi một chỉ số tổng hợp tụt giảm --- một đóng góp về phương pháp chẩn đoán có hệ thống thay vì báo cáo đơn chỉ số.

\emph{Thứ ba}, đề tài giới thiệu \textbf{audit bridge} làm phép đo điều kiện tiên quyết cho mọi đánh giá end-to-end. Phép đo này thường bị bỏ qua trong các nghiên cứu chuẩn hoá địa chỉ, dẫn đến diễn giải sai khi metric giảm sút. Việc tách audit khỏi đo chất lượng mô hình là một thực hành phương pháp luận có giá trị chuyển giao.

\emph{Thứ tư}, đề tài thiết lập \textbf{quy trình thực nghiệm phân tầng K~=~5} với bảng tổng hợp \thuat{persist} trong \schema{ath}, kết hợp cơ cấu cohort \texttt{strat-v1} bao quát bốn kiểu khó (D1--D4). Đóng góp ở chỗ chứng minh khả năng đo độ ổn định thống kê (\thuat{stability}) trên một bài toán nghiệp vụ, một thực hành chưa phổ biến trong các nghiên cứu chuẩn hoá địa chỉ trước đây.

\emph{Thứ năm}, đề tài thực hiện \textbf{nghiên cứu ablation quy mô lớn đầu tiên} cho bài toán chuẩn hóa địa chỉ Việt Nam với \VNAIABLATIONTotalSpecimens\, specimens trên \VNAIABLATIONNumConfigs\, cấu hình (Mục~\ref{sec:exp-ablation}). Khác với các nghiên cứu trước chỉ so sánh 2--3 mô hình trên tập nhỏ (\(< 1.000\) mẫu) mà không báo cáo độ lệch chuẩn, ablation này chứng minh vai trò của từng thành phần (NER, retrieval, LLM) với ý nghĩa thống kê đầy đủ và provenance có thể tái lập (git commit, seed, noise profile, platform). Kết quả cho thấy retrieval là thành phần then chốt (A3: 60{,}98\% vs A4: 8{,}46\%), trong khi LLM đóng góp \(+\VNAIABLATIONLLMContributionPp\) điểm phần trăm khi kết hợp đúng --- đây là bằng chứng định lượng đầu tiên cho kiến trúc hybrid trên miền địa chỉ Việt Nam.

\subsection{Đóng góp thực tiễn}
\label{sec:concl-practical}

\emph{Thứ nhất}, đề tài cung cấp \textbf{giải pháp khắc phục hiện tượng đứt gãy dữ liệu} do cải cách hành chính 2025. Bằng ánh xạ địa chỉ tiền cải cách sang định danh hậu cải cách qua quan hệ \macode{MERGES\_INTO}, hệ thống duy trì tính liên tục dữ liệu khách hàng trong CRM và logistics, trực tiếp giảm tỷ lệ giao hàng thất bại (\thuat{last-mile failure rate}).

\emph{Thứ hai}, kiến trúc \textbf{Waterfall Enrichment} ba lớp Cache--OSM/VietMap--Google Maps cùng chiến lược ``Fast and Free'' (toàn bộ thành phần \thuat{self-hosted} mã nguồn mở) giảm đáng kể chi phí API thương mại so với phụ thuộc \(100\%\) Google Maps. Đây là giải pháp khả thi cho doanh nghiệp Việt Nam quy mô vừa và nhỏ.

\emph{Thứ ba}, đề tài đưa lựa chọn \textbf{Qwen2.5-1.5B-Instruct với lượng tử hoá 4-bit} thành tham chiếu cụ thể về cân bằng giữa độ chính xác và hạ tầng. Cấu hình này chạy được trên một GPU tiêu dùng \(8\) GB VRAM, định vị rõ tính khả thi triển khai trong môi trường doanh nghiệp Việt Nam.

\emph{Thứ tư}, mọi mã nguồn, schema dữ liệu, script benchmark và file cấu hình đều được tổ chức để có thể tái sử dụng. Doanh nghiệp có thể bật tự động hoá theo ngưỡng ACS, giữ con người trong vòng lặp (\thuat{human-in-the-loop}) trên tập lỗi audit hoặc \thuat{stratum} khó, và mở rộng dần miền dữ liệu mà không phải tái thiết kế.

% =====================================================================
\section{Khẳng định Tính Thực tiễn của Khung Giải pháp}
\label{sec:concl-practical-affirm}

Tính thực tiễn của khung giải pháp được khẳng định trên ba căn cứ định lượng. \emph{Thứ nhất}, khung cho phép tách bạch các giai đoạn (trích thực thể, truy hồi ứng viên, suy luận cấu trúc, kiểm chứng lineage và không gian) và \emph{đo lường từng tầng riêng biệt} trên dữ liệu thật. \emph{Thứ hai}, quy mô dữ liệu được kiểm chứng đủ tin cậy về mặt thống kê: \VNAIGENAuditQueueTotal\, bản ghi hàng đợi, cohort ground truth nhiều nghìn mẫu, và đặc biệt là \VNAIABLATIONTotalSpecimens\, specimens trong nghiên cứu ablation pipeline thật. \emph{Thứ ba}, khung đã chứng minh \emph{khả năng tái lập} của thực nghiệm có kiểm soát (SUPA, ablation, log huấn luyện đầy đủ \thuat{provenance}). Trên căn cứ này, có thể khẳng định khung giải pháp \emph{vượt mức một mô hình đơn lẻ trên bảng kết quả nhỏ} (\thuat{table-running paper}) để đạt cấp độ \textbf{kiến trúc nghiên cứu có thể triển khai} (\thuat{deployable research}).

Cách tiếp cận này phù hợp triển khai trong doanh nghiệp với ba đặc tính: (a) bật tự động hoá theo ngưỡng tin cậy ACS một cách an toàn; (b) giữ con người trong vòng lặp trên tập lỗi audit hoặc \thuat{stratum} khó; (c) mở rộng dần miền dữ liệu và mô hình mà không phải tái cấu trúc khung. Đây chính là phương châm thực dụng (\thuat{pragmatic deployment}) mà các hệ thống AI hiện đại đang theo đuổi.

% =====================================================================
\section{Các Hạn chế của Nghiên cứu}
\label{sec:concl-limitations}

Đề tài chủ động công bố bốn nhóm hạn chế còn tồn tại để các vòng nghiên cứu tiếp nối có cơ sở khắc phục.

\emph{Hạn chế về dữ liệu và miền (domain).} NER được huấn luyện trên một tập công khai giới hạn \(\VNAIGENNERTrainN\,/\,\VNAIGENNEREvalN\) mẫu trong một lần đã đóng băng; khả năng khái quát hoá sang địa chỉ nội bộ doanh nghiệp, viết tắt vùng miền hoặc chuỗi từ nhận dạng ký tự quang học (OCR) cần đánh giá bổ sung. Ground truth từ Typesense và hàng đợi có thể tồn tại sự lệch phân phối (\thuat{distribution shift}) so với dữ liệu sản phẩm.

\emph{Hạn chế về đo lường chưa đủ mảnh ghép.} Mặc dù ablation pipeline thật trên \VNAIABLATIONTotalSpecimens\, specimens đã cung cấp các số liệu E2E định lượng có ý nghĩa thống kê (Mục~\ref{sec:exp-ablation}), một số mảnh ghép vẫn còn thiếu để có bức tranh ``định lượng toàn pipeline'' hoàn chỉnh trong một tài liệu duy nhất: chỉ số retrieval R@k/MRR riêng biệt, kịch bản S1--S3 trên địa chỉ thực tế, và ablation trên cohort phân tầng D1--D4 để đánh giá độ bền ở các tầng nhiễu nặng. Hạ tầng đo lường (script, schema, API) đã sẵn sàng và việc bổ sung không yêu cầu thiết kế lại.

\emph{Hạn chế phụ thuộc hạ tầng không gian.} Các endpoint Point-in-Polygon và \thuat{mismatch-report} yêu cầu PostGIS được cài đặt và \macode{mat.area\_polygon} đủ phủ. Khi thiếu, các chỉ báo không gian suy giảm xuống mức heuristic (\thuat{nearest centroid}, edge inject). Đề tài đã chuẩn bị \thuat{fallback}, song mức chính xác bị giới hạn.

\emph{Rủi ro về mô hình ngôn ngữ lớn được kiểm soát qua kiến trúc hybrid.} LLM về nguyên tắc có thể sinh chuỗi hợp ngữ pháp nhưng lệch master hoặc lệch ứng viên retrieval (hiện tượng \thuat{hallucination}). Tuy nhiên, ablation trên \VNAIABLATIONTotalSpecimens\, specimens (Mục~\ref{sec:exp-ablation}) đã định lượng được rủi ro này: cấu hình A4\_NER\_LLM (LLM không có retrieval làm khung tham chiếu) thất bại nghiêm trọng với EM chỉ \VNAIABLATIONAFourEMvTwoPct\,\% và F1 Phường giảm xuống \VNAIABLATIONAFourFOneWardPct\,\%, trong khi A1\_FULL (LLM tinh chỉnh trên top-\(k\) ứng viên retrieval) đạt EM \VNAIABLATIONAOneEMvTwoPct\,\%. Như vậy, rủi ro hallucination \emph{đã được khắc phục về mặt cấu trúc} bằng kiến trúc lai ghép, không phải bằng các kỹ thuật tinh chỉnh LLM phức tạp. Các kỹ thuật bổ sung (ràng buộc schema JSON, \thuat{constrained decoding} bám top-\(k\), RLHF trên domain) vẫn được khuyến nghị cho các vòng nâng cấp tiếp theo nhằm đẩy hiệu năng vượt ngưỡng EM \(\geq 85\,\%\) trên cohort phân tầng đầy đủ.

\emph{Hạn chế về chi phí và độ trễ.} Pipeline đầy đủ (nhúng corpus, NER, LLM) tiêu tốn đáng kể tài nguyên GPU. Đo lường thực tế trên GPU Colab T4 (Mục~\ref{sec:exp-ablation}) cho thấy A1\_FULL đạt độ trễ trung bình \VNAIABLATIONAOneLatencyMs\, ms/mẫu, cao hơn cấu hình không LLM (\VNAIABLATIONATwoLatencyMs\, ms) khoảng \VNAIABLATIONLatencyOverheadFactor\, lần --- một mức đánh đổi chấp nhận được cho các ứng dụng yêu cầu độ chính xác cao và vẫn nằm trong ngưỡng phục vụ trực tuyến \(< 200\) ms. Cần bổ sung bảng đo trên phần cứng mục tiêu triển khai cụ thể (GPU sản xuất, CPU-only fallback) cùng phân tích P95 và throughput trong chế độ batch.

% =====================================================================
\section{Hướng Phát triển}
\label{sec:concl-future}

Trên cơ sở các kết quả đã đạt được và hạn chế đã nhận diện, đề tài đề xuất sáu hướng phát triển theo trình tự tăng dần độ tự động hoá và giảm dần can thiệp thủ công, đặc biệt tập trung vào \emph{tự động hoá hoàn toàn quy trình làm giàu dữ liệu không gian} (end-to-end geospatial enrichment).

\subsection{Chuẩn hoá hạ tầng không gian}
\label{sec:concl-h1}

Đảm bảo PostGIS được kích hoạt và chỉ mục không gian (GiST) trên các cột geometry của \\ \macode{mat.area\_polygon}. Thống nhất SRID \(4326\) (WGS84) và quy trình nhập/cập nhật polygon từ OpenStreetMap hoặc nguồn ranh giới chính thống. Theo dõi độ phủ (\thuat{coverage}) theo cấp tỉnh, huyện, xã và phát hiện các vùng có dữ liệu thưa.

\subsection{Pipeline làm giàu đóng vòng (Closed Loop)}
\label{sec:concl-h2}

Sau bước \thuat{geocode} hoặc khi có toạ độ từ retrieval corpus, kích hoạt chuỗi tự động: gán đơn vị hành chính bằng Point-in-Polygon; nếu \thuat{mismatch} so với trường khai báo, kích hoạt Edge Injection hoặc đưa vào hàng đợi xử lý ngoại lệ; ghi nhận nguồn và độ tin cậy vào metadata bản ghi. Đây là sự kết nối trực tiếp giữa Module Geospatial (\ref{sec:geospatial}) với Module AI Pipeline (\ref{sec:ai-pipeline}) đã thiết kế.

\subsection{Tự động hoá thu thập OSM}
\label{sec:concl-h3}

Lên lịch job thu thập OpenStreetMap theo tỉnh với ngưỡng đầy đủ dữ liệu đường và POI; đồng bộ \thuat{incremental} chỉ với các thay đổi khi API Overpass và lưu trữ cho phép; gắn SLA cụ thể về độ trễ và kích thước job. Phương án này giảm tải hạ tầng so với tải lại toàn bộ dữ liệu OSM định kỳ.

\subsection{Liên kết với chuẩn hoá địa chỉ}
\label{sec:concl-h4}

Áp dụng cùng một ``hợp đồng dữ liệu'' (\thuat{data contract}) lineage queue--master đã được kiểm định bởi audit bridge. Ưu tiên \thuat{join} theo lineage cho mọi phân tích chất lượng; giảm dần tỷ lệ \thuat{fail} G3 (\VNAIGENAuditGThreeFail\, dòng) bằng quy tắc sửa hoặc \thuat{re-sync} master định kỳ.

\subsection{Đánh giá và Giám sát Liên tục}
\label{sec:concl-h5}

Bổ sung báo cáo định kỳ cho retrieval và E2E; tích hợp dashboard theo \thuat{stratum} (tỉnh, độ dài chuỗi, epoch hành chính); thiết lập cảnh báo khi Gate~B hoặc các cổng audit suy giảm sau migration dữ liệu hoặc nâng cấp mô hình. Trên cơ sở kết quả ablation đã xác nhận A1\_FULL là kiến trúc tối ưu (Mục~\ref{sec:exp-ablation}), hai mục tiêu đo lường ưu tiên gồm: (i) chạy A1\_FULL trên cohort phân tầng K~=~5 với \(n = 2.000\) mỗi lần để khẳng định độ ổn định và độ bền trên các tầng D1--D4 (đặc biệt D2 nhiễu cao và D3 lưỡng thời); (ii) chạy SUPA Final với \(N = 10.000\) trên A1\_FULL để xác nhận mức EM ổn định ở quy mô gấp đôi ablation hiện tại. Đây là chuyển từ ``chạy benchmark một lần'' sang ``vận hành benchmark như dịch vụ'' (\thuat{benchmarking as a service}).

\subsection{Tự động hoá pipeline huấn luyện liên tục}
\label{sec:concl-h5bis}

Hệ thống đã hiện thực cơ chế huấn luyện liên tục PhoBERT NER từ các mẫu PreLabeler đã được con người xác nhận đạt chuẩn (\macode{ai.prelabeler\_testcases} với \macode{test\_result->>'passed' = true}). Endpoint \macode{POST /api/training/trigger-ner-from-prelabeler} cho phép kích hoạt huấn luyện khi số lượng mẫu mới đủ ngưỡng (mặc định 50). Hướng phát triển là tự động hoá hoàn toàn vòng lặp: \emph{người dùng gán nhãn} \(\to\) \emph{validation tự động} \(\to\) \emph{huấn luyện} \(\to\) \emph{đánh giá trên tập kiểm định cố định} \(\to\) \emph{deploy nếu F1 cải thiện}; mỗi checkpoint mới ghi đồng thời vào \macode{ath.training\_history} với \macode{notes} dạng \texttt{prelabeler\_cases=...;labelstudio\_merged=...;eval\_loss=...} và git commit hash, đảm bảo provenance đầy đủ và khả năng rollback an toàn.

\subsection{Nghiên cứu Nâng cao}
\label{sec:concl-h6}

Bốn hướng nghiên cứu nâng cao đã được nhận diện. \emph{Một}, học có cấu trúc (\thuat{structured prediction}) kết hợp ràng buộc đồ thị hành chính trên \macode{mat.unit\_edge}. \emph{Hai}, học từ phản hồi con người (\thuat{learning from human feedback}) trên tập ngoại lệ \thuat{mismatch}, cho phép mô hình tự cải thiện theo phản hồi vận hành. \emph{Ba}, mô hình sinh có ràng buộc (\thuat{constrained decoding}) khớp tập ứng viên retrieval, khắc phục \thuat{hallucination} của LLM. \emph{Bốn}, học liên kết (\thuat{federated learning}) cho phép nhiều đơn vị logistics cùng huấn luyện mô hình mà không cần chia sẻ dữ liệu khách hàng nhạy cảm --- một hướng đặc biệt phù hợp với Luật Bảo vệ Dữ liệu Cá nhân của Việt Nam.

% =====================================================================
\section{Lời kết}
\label{sec:concl-final}

Đề tài \emph{Xây dựng khung giải pháp làm giàu và chuẩn hoá dữ liệu địa chỉ Việt Nam sử dụng tiếp cận đa nguồn và thuật toán hình học không gian trong bối cảnh sắp xếp đơn vị hành chính toàn quốc 2025} đã hoàn thành mục tiêu xây dựng khung giải pháp VNAI với ba kết luận khoa học chính:

\textbf{Một, kiến trúc hybrid (NER + retrieval + LLM) là cần thiết và đủ.} Nghiên cứu ablation trên \VNAIABLATIONTotalSpecimens\, specimens chứng minh: (i)~retrieval là thành phần then chốt không thể bỏ qua (A3\_MGTE\_ONLY: 60{,}98\% vs A4\_NER\_LLM: 8{,}46\%); (ii)~LLM đóng góp \(+\VNAIABLATIONLLMContributionPp\) điểm phần trăm khi kết hợp đúng (A1\_FULL: 66{,}58\% vs A2/A3: 60{,}98\%); (iii)~TF-IDF và mGTE tương đương trên cohort phổ thông. Đây là bằng chứng định lượng đầu tiên cho kiến trúc hybrid trên bài toán địa chỉ Việt Nam với quy mô đủ lớn để có ý nghĩa thống kê.

\textbf{Hai, mô hình SCD Type 2 + unit\_edge graph giải quyết được vấn đề lưỡng thời hành chính.} Khác với các nghiên cứu trước coi dữ liệu hành chính là tĩnh, VNAI cho phép xử lý đồng thời địa chỉ tiền và hậu cải cách 2025 mà không gặp \thuat{catastrophic forgetting}. Đồ thị quan hệ \macode{mat.unit\_edge} với bốn kiểu quan hệ (MERGES\_INTO, SPLIT\_FROM, RENAMES\_TO, BOUNDARY\_ADJUSTED) kết hợp với cơ chế SCD Type 2 tạo nền tảng cho \emph{Temporal-Aware Address Standardization} --- đây là đóng góp lý thuyết chính của đề tài.

\textbf{Ba, phương pháp luận SUPA-Bench + Audit Bridge + Ablation quy mô lớn có thể chuyển giao.} Khung đánh giá với provenance đầy đủ (git\_commit, seed, noise\_profile), tách bạch chỉ số P-*/S-*, và ablation study \(N = 25.000\) là đóng góp phương pháp có giá trị cho các nghiên cứu tương tự. Chiến lược thực nghiệm phân tầng (NER \(\to\) Audit \(\to\) Oracle \(\to\) K~=~5 \(\to\) Ablation) đảm bảo mỗi kết luận đều có bằng chứng định lượng vững chắc.

Khung giải pháp VNAI không tham vọng thay thế luật phân cấp hành chính do Nhà nước ban hành; thay vào đó, đề tài hướng tới mục tiêu \emph{kéo suy luận của AI vào đúng bằng chứng trên master hành chính} và \emph{báo cáo các chỉ tiêu theo giao thức đã khoá}, để mọi cải tiến đều có thể đo lường và đối chứng. Cấu hình tối ưu A1\_FULL (NER + mGTE + LLM) đã được xác lập với EM@v2 \(= \VNAIABLATIONAOneEMvTwoPct\,\%\), vượt nhiều ngưỡng kỳ vọng nội bộ.

Các mảnh ghép định lượng còn lại --- ablation phân tầng K~=~5 trên cohort D1--D4, SUPA Final \(N = 10.000\) trên pipeline A1\_FULL, chỉ số retrieval R@k riêng biệt, và bảng latency P95/throughput trên hạ tầng production --- là khoảng trống công khai cho các vòng nghiên cứu tiếp nối. Kính mong nhận được góp ý chuyên môn từ Hội đồng để khung giải pháp ngày càng hoàn thiện và thực sự đóng góp cho hạ tầng dữ liệu địa chỉ bền vững phục vụ nền kinh tế số Việt Nam.
